\documentclass[12pt]{article}
\usepackage[T1]{fontenc}
\usepackage[utf8]{inputenc}
\usepackage{lmodern}
\usepackage{textcomp}
\usepackage{enumitem}
\usepackage{geometry}
\geometry{margin=1in}
\begin{document}
\begin{center}
Answer Key - Physics - Graduate Level Assessment 23 \
\small (Answers are ordered exactly as in the question paper.)
\end{center}

\section*{Multiple Choice}
\begin{enumerate}[label=\arabic*.]
\item \textbf{Answer:} B
\\ \textit{Options:} A) Refraction; B) de-excitation; C) Chemiluminescence; D) Thermal radiation; E) Phosphorescence
\\ \textit{Note:} The red glow in the neon tube is produced by de-excitation of neon gas atoms, which emit light as they return to a lower energy state after being energized by an electric current.
\item \textbf{Answer:} A
\\ \textit{Options:} A) 2.5 amp; B) 0.5 amp; C) 2.8 amp; D) 3.0 amp; E) 1.0 amp
\\ \textit{Note:} The correct answer is 2.5 amps because the total impedance of the circuit, calculated using the formula Z = √(R² + X\_L²), results in 68.31 ohms, leading to a current of I = V/Z = 130 V/68.31Ω, which approximately equals 2.5 amps.
\item \textbf{Answer:} E
\\ \textit{Options:} A) 700 nm; B) 450 nm; C) 750 nm; D) 300 nm; E) 357 nm
\\ \textit{Note:} The correct answer is 357 nm because Rayleigh scattering is inversely proportional to the fourth power of the wavelength, so a decrease in wavelength from 635 nm to 357 nm results in an increase in scattering by a factor of 10.
\item \textbf{Answer:} A
\\ \textit{Options:} A) 3.2 \ensuremath{\times} $10^10$ cm/sec; B) 3.0 \ensuremath{\times} $10^10$ cm/sec; C) 2.8 \ensuremath{\times} $10^10$ cm/sec; D) 2.4 \ensuremath{\times} $10^10$ cm/sec; E) 1.8 \ensuremath{\times} $10^10$ cm/sec
\\ \textit{Note:} The correct answer, 3.2 \ensuremath{\times} 10^\{10\} cm/sec, is derived from the relativistic mass-energy relation which states that the relativistic mass of a particle increases with its velocity and is given by the equation m = m\_0 / sqrt(1 - ($v^2$/$c^2$)), where the mass becomes twice the rest mass (m = 2 m\_0) at that specific velocity.
\item \textbf{Answer:} A
\\ \textit{Options:} A) salinity of the water.; B) proximity to the equator.; C) presence of fish in the pond.; D) type of pond vegetation.; E) acceleration due to gravity.
\\ \textit{Note:} The pressure at the bottom of a pond is determined solely by the depth of the water column and the density of the water, which is primarily influenced by gravitational force, rather than the salinity of the water.
\end{enumerate}

\section*{True / False}
\begin{enumerate}[label=\arabic*.]
\setcounter{enumi}{5}
\item \textbf{Answer:} False
\\ \textit{Options:} A) True; B) False
\\ \textit{Note:} The Mars Exploration Rover Spirit was actually tilted towards the south to optimize solar power generation during winter and to minimize exposure to dust accumulation on its solar panels.
\item \textbf{Answer:} False
\\ \textit{Options:} A) True; B) False
\\ \textit{Note:} The statement is false because the reflectance at a water-glass boundary depends on the refractive indices of the materials, and the given reflectance values for crown glass and flint glass are inaccurately specified; for normal incidence, the reflectance for crown glass is approximately 2.1\% and for flint glass it is around 4.0\%.
\item \textbf{Answer:} True
\\ \textit{Options:} A) True; B) False
\\ \textit{Note:} The frequency of the next higher harmonic for a closed organ pipe is calculated using the formula for harmonics of closed pipes, which states that the harmonics are odd multiples of the fundamental frequency, making the first harmonic (fundamental) 131 Hz and the next harmonic (third) 3 times the fundamental, resulting in 393 Hz, indicating that the statement is actually false, as 196 Hz does not correspond to any harmonic of 131 Hz.
\item \textbf{Answer:} False
\\ \textit{Options:} A) True; B) False
\\ \textit{Note:} The statement is false because an object can only have zero acceleration when the net external forces are balanced to zero, which means the individual forces must be equal in magnitude and opposite in direction, not just perfectly balanced in some other sense.
\item \textbf{Answer:} True
\\ \textit{Options:} A) True; B) False
\\ \textit{Note:} The impedance of an inductor is given by the formula Z = jωL, where Z is the impedance, j is the imaginary unit, ω is the angular frequency, and L is the inductance; thus, for a 1-henry inductor, the impedance at 100 rad/sec is 100 j ohms, at 1000 rad/sec is 1000 j ohms, and at 10,000 rad/sec is 10,000 j ohms, making the statement about the magnitude of the impedance at these frequencies correct.
\end{enumerate}

\section*{Long Form Response}
\begin{enumerate}[label=\arabic*.]
\setcounter{enumi}{10}
\item \textbf{Answer:} To calculate the height of the tree, we can apply the principles of similar triangles and the geometry of reflection. Given that the observer is 100 m away from the tree and the mirror is positioned 30 cm (0.3 m) from the observer, the total distance from the tree to the mirror is 100 m - 0.3 m = 99.7 m. The height of the image reflected in the mirror is equal to the height of the tree, and the mirror's height is 5 cm (0.05 m). Using the ratio of the distances, we find that the height of the tree can be calculated as follows: (height of the mirror / distance from observer to mirror) = (height of the tree / distance from tree to mirror). Solving this proportion results in a tree height of approximately 16.67 m. This scenario illustrates the principles of optics, specifically the behavior of light and reflection in forming images.
\\ \textit{Note:} correct because it utilizes the principles of similar triangles and reflection geometry to establish the relationship between the distances of the observer, mirror, and tree to accurately calculate the tree's height based on the height of the mirror and the distances involved.
\item \textbf{Answer:} Operating a heat engine, such as a jet engine, at high temperatures significantly enhances its efficiency due to the principles outlined in the second law of thermodynamics and the Carnot efficiency equation. Higher temperatures increase the temperature difference between the heat source and the heat sink, which directly boosts the maximum theoretical efficiency of the engine, as described by the formula η = 1 - (Tc/Th), where Tc is the cold reservoir temperature and Th is the hot reservoir temperature. Additionally, high operating temperatures lead to improved thermodynamic properties of the working fluid, such as higher energy content and lower viscosity, which facilitate greater energy conversion and reduced frictional losses. However, these advantages must be balanced against material limitations and the potential for increased thermal stress, which can lead to mechanical failure if not carefully managed. Thus, while high-temperature operation can dramatically enhance efficiency, it necessitates advanced materials and engineering solutions to ensure reliability and performance.
\\ \textit{Note:} The answer correctly identifies that operating a heat engine at higher temperatures increases the efficiency by maximizing the temperature difference between the heat source and sink, as dictated by the second law of thermodynamics and illustrated by the Carnot efficiency equation.
\end{enumerate}

\vfill
\noindent\textit{End of Answer Key}
\end{document}